


\newpage

%\footline={\hss\folio\hss}
%\pageno=1

\dspace
\subhead
5. Le groupe fondamental motivique du
compl\'ement, dans $\dbP^1$, de $0$, $\infty$,
et $\mu_N$
\endsubhead

\subhead
5.1
\endsubhead
Soient $N$ un entier $\Ge 1$,
$k$ engendr\'e sur $\dbQ$ par une racine
primitive $N^{\ieme}$ de l'unit\'e
et $\scrO$ l'anneau de ses entiers.
Il est commode de ne pas supposer que $k$ soit
le sous-corps $\dbQ(\exp(2\pi i/N))$ de
$\dbC$.
Notons $A_k$, ou simplement $A$ la droite
affine type sur $k$ et posons
$A^*=A-\{0\}=\Spec\,k[u,u^{-1}]$.
Soit $X$ le compl\'ement dans $A^*$ de
l'ensemble $\mu_N:=\mu_N(k)$ des racines
$N^{\iemes}$ de l'unit\'e.
S'il y a ambigu{\umi}t\'e sur $N$, on \'ecrira
plut\^{o}t $k_N$, $\scrO_N$ et $X_N$.

Le groupe dih\'edral $\dbZ/2\ltimes\mu_N$ agit
sur $X\Subset \dbP_k^1$: le g\'en\'erateur de
$\dbZ/2$ agit par $u\mapsto u^{-1}$, et
$\xi\in\mu_N$ agit par $u\mapsto\xi u$.
Pour $N=1,2$ ou $4$, le groupe des
projectivit\'es transformant $X$ en
lui-m\^{e}me est plus grand.
Pour $N=1$, il s'identifie au groupe des
permutations de $\{0,1,\infty\}$.
Pour $N=2$, $(0,\infty,1,-1)$ est un quaterne
harmonique, et on obtient le groupe des
automorphismes d'un carr\'e de sommets
cons\'ecutifs $0$, $1$, $\infty$, $-1$.
Pour $N=4$, les points $0$, $1$, $i$, $-1$, $-i$,
$\infty$ forment sur la sph\`ere de Riemann les
sommets d'un octa\`edre, et on obtient le
groupe des d\'eplacements de cet octa\`edre.

\subhead
5.2
\endsubhead
Si $x$ et $y$ sont deux points-base (4.3) de
$X$, l'espace de chemins motivique $P_{y,x}$
est un $\MT(k)$-sch\'ema affine (4.4).
La r\'ealisation $\omega$ de $P_{y,x}$ est
ind\'ependante de $x$ et $y$.
C'est un $\dbQ$-sch\'ema en groupe pro-unipotent, 
et la composition des chemins est
la loi de groupe.
Autre fa\c{c}on de dire: en r\'ealisation
$\omega$, il y a un chemin canonique de $x$
\`a $y$, et le compos\'e des chemins
canoniques de $x$ \`a $y$ et de $y$ \`a 
$z$ est le chemin canonique de $x$ \`a $z$.
Bien s\^{u}r, cette trivialisation du
$\omega(\pi_1(X,x))$-torseur $\omega(P_{y,x})$
n'est en g\'en\'eral pas motivique, i.e. ne
provient pas d'une trivialisation du
$\pi_1(X,x)$-torseur $P_{y,x}$.
L'action sur $\omega(P_{y,x})$
du groupe des automorphismes du foncteur fibre
$\omega$ d\'epend de $x$ et $y$.

\subhead
5.3
\endsubhead
Pour $z$ dans $\dbP_k^1-\{\infty\}$ et
$\lambda$ dans $k$ nous noterons
$\lambda_z$ le vecteur tangent en $z$ de
coordonn\'ee
$du(\lambda_z)=\lambda$.
Nous noterons $\lambda_\infty$ le vecteur
tangent en $\infty$ de coordonn\'ee
$du^{-1}(\lambda_\infty)=\lambda$.

Nous prendrons comme points-base les
$\lambda_z$ pour $z=0$, $\infty$ ou dans
$\mu_N$ et $\lambda$ une racine de l'unit\'e
dans $k$.
Ce syst\`eme de points-base est stable sous
l'action du groupe dih\'edral
$\dbZ/2\ltimes\mu_N$.
Une application de 1.7 montre comme en 4.17,
4.18 que les $P_{y,x}$ correspondants sont non
ramifi\'es en dehors de 
l'ensemble des places de $k$ au-dessus
d'un nombre premier divisant $N$: ce sont
des $\MT(\scrO[1/N])$-sch\'emas.
Nous d\'ecrivons en 5.4 et 5.5 des structures
sur le syst\`eme des $P_{y,x}$.
Pour toutes ces structures, il suffit pour les
construire de les construire dans
$\scrR^{H+\ell}$, ce qui est fait dans
Deligne (1989), et d'appliquer 2.14.


\subhead
5.4
\endsubhead
Soient $T$ une droite (sch\'ema vectoriel de
dimension un) et $T^*=T-\{0\}$.
Le $\pi_1$ motivique de $T^*$ est ab\'elien,
ind\'ependant du point-base.
C'est $\dbQ(1)$.
Pour $x\in T^*(k)$, l'application
$\lambda\mapsto\lambda x$ est un isomorphisme
de $A^*$ avec $T^*$.
Pour $y\in T^*(k)$, il induit un isomorphisme
de $\dbQ(1)$-torseurs de
$P_{y/x,1}(A^*)$ avec $P_{y,x}(T^*)$.
Le $\dbQ(1)$ torseur $P_{t,1}(A^*)$
est le {\it torseur de Kummer} $K(t)$, et la
composition des chemins fournit un
isomorphisme de torseurs $K(tu)=K(t)+K(u)$.

En particulier, si $\alpha\in k^*$ est une
racine $n^{\ieme}$ de l'unit\'e, elle fournit
une trivialisation de $nK(\alpha)$ et, puisque la
multiplication par $n$ est un automorphisme de
$\dbQ(1)$, par division, une trivialisation de
$K(\alpha)$.

Pour $x\in T^*(k)$ et $\alpha$ une racine de
l'unit\'e, le $\dbQ(1)$-torseur $P_{\alpha
x,x}$ est ainsi trivialis\'e: on dispose d'un
chemin motivique de $x$ \`a $\alpha x$.

Prenons pour $T$ l'espace tangent \`a
$\dbP^1$ en $0$, $\infty$ ou en
$\zeta\in\mu_N$.
Le groupo{\umi}de des $P_{y,x}(T^*)$, pour
$x,y\in T^*(k)$, s'envoie dans le 
groupo{\umi}de
des $P_{y,x}(X)$.
Pour $x\in T^*(k)$ et $\alpha$ une racine de
l'unit\'e on dispose donc d'un chemin
motivique de $x$ \`a $\alpha x$ dans $X$: une
trivialisation du $\pi_1(X,x)$-torseur
$P_{\alpha x,x}(X)$.

Pour le syst\`eme de points base $\lambda_z$
de 5.3, $P_{\lambda_z,\mu_t}$ ne d\'epend donc
que de $z$ et $t$, et on \'ecrira simplement
$P_{z,t}$ pour $P_{\lambda_z,\mu_t}$.
Quand d'autres points-base \`a l'infini auront \`a
\^{e}tre utilis\'es, nous retournerons \`a la
notation $\lambda_z$.


\subhead
5.5
\endsubhead
On dispose sur le syst\`eme des
$\MT(\scrO[1/N])$-sch\'emas $P_{y,x}$ ($x,y\in
S:=\{0,\infty\}\Cup \mu_N$) des structures
suivantes.

\medskip\noindent
{\sans (A)}\enspace
La composition des chemins.

\smallskip\noindent
{\sans (B)}\enspace
Pour chaque $x\in S$, un 
morphisme de sch\'emas en groupe, la {\it
monodromie locale autour de $x$}:
$$
\dbQ(1)\to P_{x,x}.
$$

\smallskip\noindent
{\sans (C)}\enspace
une \'equivariance sous le groupe dih\'edral
$\dbZ/2\ltimes\mu_N$ de 5.1.

\subhead
5.6
\endsubhead
Si on applique un foncteur fibre $F$, \`a
valeurs dans les espaces vectoriels sur un
corps $K$, \`a $\dbQ(1)$, aux $P_{y,x}$ et 
aux structures 5.5, on obtient une structure
de l'esp\`ece suivante:

\medskip\noindent
(1)\enspace
Un espace vectoriel de dimension un $K(1)$.

\noindent
(2)\enspace
Pour chaque $x,y\in S$, un sch\'ema $P_{y,x}$
sur $K$.

\noindent
$(3)_{\text{A}}$\enspace
Un syst\`eme de morphismes de sch\'emas
$$
P_{z,y}\times P_{y,x}\to P_{z,x}\tag5.6.1
$$
faisant des $P_{y,x}$ un groupo{\umi}de.
Les sch\'emas en groupe $P_{x,x}$ sont
pro-unipotents.

\noindent
$(3)_{\text{B}}$\enspace
Pour chaque $x\in S$, un homomorphisme
$$
\text{(groupe additif }
K(1))\to P_{x,x}.\tag5.6.2
$$
L'alg\`ebre de Lie de $P_{x,x}$ est une
pro-alg\`ebre de Lie.
Parce que $P_{x,x}$ est pro-unipotent, la
donn\'ee (5.6.2) \'equivaut \`a celle de
$$
K(1)\to \Lie\,P_{x,x}.\tag5.6.3
$$

\noindent
$(3)_{\text{C}}$\enspace
Une $\dbZ/2\ltimes\mu_N$ \'equivariance.

\subhead
5.7
\endsubhead
Il nous sera utile d'affaiblir une structure
5.6 en restreignant les points-base \`a
\^{e}tre $0$ ou dans $\mu_N$.
Plus pr\'ecis\'ement, \`a une structure 5.6,
on peut attacher une structure affaiblie du
type suivant.

\medskip\noindent
(1)\enspace
Comme en 5.6 (1).

\noindent
(2)\enspace
Un sch\'ema en groupe pro-unipotent $P_{0,0}$
et, pour $\zeta\in\mu_N$, un $P_{00}$-torseur
$P_{\zeta,0}$.
On note $P_{\zeta\zeta}$ le tordu de $P_{00}$
par ce torseur, i.e. le sch\'ema en groupe des
automorphismes du $P_{0,0}$-torseur
$P_{\zeta,0}$.

\noindent
$(3)_{\text{B}}$\enspace
Pour $x=0$ ou $x\in\mu_N$, un morphisme
(5.6.2) ou, ce qui revient au m\^{e}me,
(5.6.3).

\noindent
$(3)_{\text{C}}$\enspace
Une $\mu_N$-\'equivariance.

\remark{\bf Remarques} (i)
La donn\'ee d'une structure 5.7 \'equivaut \`a
celle de (1) et de

\noindent
$(2')$\enspace
Un sch\'ema en groupe pro-unipotent $P_{00}$
et un $P_{00}$-torseur $P_{10}$.
Comme en (2), $P_{10}$ d\'efinit par torsion
$P_{11}$.

\noindent
$(3')_{\text{B}}$\enspace
Comme en $(3)_{\text{B}}$, pour $x=0$ ou $1$.

\noindent
$(3')_{\text{C}}$\enspace
Une action sur $P_{00}$ de $\mu_N$, respectant
la monodromie locale autour de $0$.

On r\'ecup\`ere (2) en d\'efinissant
$P_{\zeta,0}$ comme une copie de $P_{1,0}$,
l'\'equivariance $(3)_{\text{C}}\,:\zeta: P_{1,0}\to
P_{\zeta,0}$ \'etant l'identit\'e de $P_{1,0}$
sur cette copie.

\noindent
(ii)\enspace
Cette structure est un peu plus pr\'ecise que
la donn\'ee de $P_{00}$, de l'action de
$\mu_N$, de $(3)_{\text{B}}$ pour $x=0$ et de
la classe de conjugaison ``monodromie locale
autour de $1$'' de morphismes $\dbQ(1)\to
P_{1,1}$: cette classe de conjugaison est un
espace homog\`ene, qu'on peut identifier au
quotient $P_{0,1}/\dbQ(1)$.
\endremark

\subhead
5.8
\endsubhead
Prenons pour foncteur fibre le foncteur fibre
$\omega$, et d\'ecrivons la structure
d'esp\`ece (5.6) obtenue.
Puisque $\omega$ est \`a valeurs dans les
$\dbQ$-espaces vectoriels gradu\'es, les
alg\`ebres affines des $\omega(P_{y,x})$ seront
des $\dbQ$-alg\`ebres gradu\'ees.

Soit $\scrL$ l'alg\`ebre de Lie engendr\'ee
par des g\'en\'erateurs $e_x$, 
$x\in\{0,\infty\}\Cup\mu_N$, soumis \`a la seule
relation que $\sum e_x=0$.
En cas d'ambigu{\umi}t\'e, on l'\'ecrira plus
pr\'ecisement $\scrL_N$.
On la gradue en donnant aux $e_x$ le degr\'e un.
Elle est librement engendr\'ee par les $e_x$,
$x\not=\infty$.

Soit $\prod$ le groupe pro-unipotent
$$
\prod:=\lim\,\exp(\scrL/\text{degr\'e $\Ge
n$)}.
$$
On munit son alg\`ebre affine de la graduation
d\'eduite de celle de $\scrL$.
L'alg\`ebre de Lie de $\prod$ est le
compl\'et\'e $\scrL^{\uphat}
:=\prod\limits_{n}\scrL_n$ de $\scrL$.

Avec ces notations, la structure d'esp\`ece
5.6, et les graduations, sont comme suit.

\medskip\noindent
$(1)^\omega$\enspace
L'espace vectoriel $\dbQ$, en degr\'e un.

\noindent
$(2)^\omega$\enspace
Pour chaque $x$, $y$, une copie
$\prod\nolimits_{y,x}$ de $\prod$.

\noindent
$(3)_{\text{A}}^\omega$\enspace
Morphismes (5.6.1): la loi de groupe de
$\prod$.

\noindent
$(3)_{\text{B}}^\omega$\enspace
Le morphisme (5.6.3) est
$$
\dbQ\to\scrL^{\uphat}\colon\, 1\longmapsto e_x
$$

\noindent
$(3)_{\text{C}}^\omega$\enspace
Pour $\sigma\in\dbZ/2\ltimes\mu_N$, agissant
sur $\{0,\infty\}\Cup\mu_N$ comme en 5.1, $\sigma$
envoie $\prod\nolimits_{y,x}$ sur
$\prod\nolimits_{\sigma y,\sigma x}$.
C'est le morphisme de groupes
$[\sigma]\colon\,\prod\to\prod$
qui induit sur l'alg\`ebre de Lie $e_x\mapsto
e_{\sigma x}$.

D'apr\`es 2.10, la m\^{e}me description vaut,
sur $k$, pour le foncteur fibre $\DR$.

Le groupe $G_\omega\left<\MT(\scrO[1/N])\right>$ 
(2.1) agit sur cette structure d'esp\`ece 5.8.
Nous nous proposons d'\'etudier cette action.


\subhead
5.9
\endsubhead
La structure 5.8 d'esp\`ece 5.6 fournit par
oubli une structure d'esp\`ece 5.7.
Pr\'esent\'ee comme en 5.7 Remarque (i), c'est
la suivante.

\medskip\noindent
$(1)^\omega$\enspace
L'espace vectoriel $\dbQ$.

\noindent
$(2)^\omega$\enspace
Une copie $\prod\nolimits_{0,0}$ de $\prod$ et
le $\prod\nolimits_{0,0}$ torseur trivial
$\prod$, not\'e $\prod\nolimits_{1,0}$.
Le tordu de $\prod\nolimits_{0,0}
$ par ce torseur est une nouvelle copie de
$\prod$, not\'ee $\prod\nolimits_{1,1}$.

\noindent
$(3)_{\text{B}}^\omega$\enspace
Les morphismes (5.6.3), pour $x=0$ ou $1$ sont
$\dbQ\to\scrL^{\uphat}\colon\, 1\to e_x$.

\noindent
$(3)_{\text{C}}^\omega$\enspace
L'action de $\mu_N$ sur
$\prod\nolimits_{0,0}$, donn\'ee sur
l'alg\`ebre de Lie $\scrL^{\uphat}$  par
$e_x\mapsto e_{\xi x}$.


\subhead
5.10
\endsubhead
Soit $H_\omega$ le sch\'ema en groupe des
automorphismes de la structure 5.9 d'esp\`ece
5.7.
L'action de $H_\omega$ sur l'espace vectoriel
de dimension un $(1)^\omega$ est un morphisme
$$
H_\omega\to \dbG_m.\tag5.10.1
$$
Soit $V_\omega$ son noyau.
Les graduations fournissent un morphisme de
$\dbG_m$ dans $H_\omega$.
Puisque $(1)^\omega$ est en degr\'e un, c'est
une section de (5.10.1).
Elle fait de $H_\omega$ un produit
semi-direct:
$$
H_\omega=\dbG_m\ltimes V_\omega.\tag5.10.2
$$

L'action de $G_\omega$ sur la structure 5.9 se
factorise par $H_\omega$, puisque $H_\omega$ a
\'et\'e d\'efini comme le groupe
d'automorphismes de r\'ealisations $\omega$
respectant des structures motiviques.
Comparant \`a 2.1, on v\'erifie que le
morphisme
$$
G_\omega=\dbG_m\ltimes U_\omega\to H_\omega
=\dbG_m\ltimes V_\omega\tag5.10.3
$$
est compatible aux d\'ecompositions en
produits semi-directs.

\proclaim{Proposition 5.11}
Notons $1_{y,x}$ l'\'el\'ement neutre de la
copie $\prod\nolimits_{y,x}$ de $\prod$.
Avec les notations de 5.9 et 5.10, le
morphisme $v\mapsto v(1_{1,0})$ est un
isomorphisme de sch\'emas de $V_\omega$ avec
$\prod\nolimits_{1,0}$.
\endproclaim

Rappelons que pour $\xi\in\mu_N$, 
$[\xi]$ est l'automorphisme de $\prod$
qui induit sur l'alg\`ebre de Lie
$e_x\mapsto e_{\xi x}$.
On notera $\prod\nolimits_{0,1}$ le
$\left(\prod\nolimits_{1,1},
\prod\nolimits_{0,0}\right)$-bitorseur inverse
du $\left(\prod\nolimits_{0,0},
\prod\nolimits_{1,1}\right)$-bitorseur
$\prod\nolimits_{1,0}$.
C'est une nouvelle copie $\prod$.
Pour $x,y,z\in\{0,1\}$, on dispose d'une
``composition des chemins''
$\prod\nolimits_{z,y}\times\prod\nolimits_{y,x}\to
\prod\nolimits_{z,x}$, simplement donn\'ee par
la loi de groupe de $\prod$.

\demo{Preuve}
Soit $v$ dans $V_\omega$ et posons
$a:=v(1_{1,0})$.
\enddemo

Puisque $V_\omega$ agit trivialement sur
$(1)^\omega$ et respecte
$(3)_{\text{B}}^\omega$, on a dans
$\Lie\,\prod\nolimits_{00}$, identifi\'e \`a
$\scrL^{\uphat}$
$$
v(e_0)=e_0\qquad\qquad
\text{(dans $\Lie\,{\tsize \prod\nolimits_{0,0})}$}
\tag5.11.1
$$
et dans $\Lie\,\prod\nolimits_{1,1}$,
\'egalement identifi\'e \`a $\scrL^{\uphat}$
$$
v(e_1)=e_1\qquad\qquad\text{(dans
$\Lie\,{\tsize \prod\nolimits_{1,1}}$)}.
$$
Transformant par $v$ l'identit\'e
$$
\left(\exp(e_1)\text{ dans }{\tsize
\prod\nolimits_{0,0}}\right)=1_{0,1} \cdot \left(\exp(e_1)
\text{ dans }{\tsize\prod\nolimits_{11}}\right)\cdot 
1_{1,0},
$$
on obtient
$$
v(\exp(e_1)\text{ dans }{\tsize \prod\nolimits_{0,0}})
=a^{-1}\exp(e_1)a,
$$
qui \'equivaut \`a
$$
v(e_1\text{ dans }\Lie\,
{\tsize\prod\nolimits_{0,0}})=
\ad_{a^{-1}}(e_1).\tag5.11.2
$$

Par \'equivariance, (5.11.2) implique que pour
$\zeta\in\mu_N$,
$$
v(e_\zeta\text{ dans }
\Lie\,{\tsize
\prod\nolimits_{0,0}})=\ad_{[\zeta]a^{-1}}
(e_\zeta).\tag5.11.3
$$

Pour $a$ dans $\prod$, notons
$\left<a\right>_0$ l'automorphisme de $\prod$
induisant sur l'alg\`ebre de
$\Lie\,\scrL^{\uphat}$ l'automorphisme, encore
not\'e $\left<a\right>_0$:
$$
\left<a\right>_0\colon\, e_0\longmapsto
  e_0,\quad
e_\zeta\longmapsto
\ad_{[\zeta]a^{-1}}(e_\zeta).
\tag5.11.4
$$
Notons $\left<a\right>_{1,0}$ l'automorphisme du
sch\'ema $\prod$
$$
\left<a\right>_{1,0}\colon\, g\longmapsto 
a \cdot \left<a\right>_0(g).\tag5.11.5
$$
Les formules (5.11.1) et (5.11.3) montrent que
$v$ agit par $\left<a\right>_0$ sur
$\prod\nolimits_{0,0}$. 
Sur le $\prod\nolimits_{0,0}$-torseur
$\prod\nolimits_{1,0}$, $v$ agit par
$\left<a\right>_{1,0}$: transformer par $v$
l'identit\'e ($g$ dans $\prod\nolimits_{1,0}$)
$=1_{1,0} \cdot (g\text{ dans
}\prod\nolimits_{0,0})$.

R\'eciproquement, quelque soit $a$ dans $\prod$,
ces formules d\'efinissent un automorphisme
$\left<a\right>$ de la structure 5.9
d'esp\`ece 5.7.

\subhead
5.12
\endsubhead
Si on transporte \`a la copie
$\prod\nolimits_{1,0}$ de $\prod$ la loi de
groupe de $V_\omega$, on obtient une nouvelle
loi de groupe $\circ$ sur le sch\'ema $\prod$.
Elle est donn\'ee par
$$
a \circ b=a \cdot \left<a\right>_0(b).\tag5.12.1
$$
En effet, $\left<b\right>_{1,0}$ envoie
$1_{1,0}\in\prod\nolimits_{1,0}$ sur $b$ et
$\left<a\right>_{1,0}$ envoie
$b=1_{1,0} \cdot b\in\prod\nolimits_{1,0}$ sur
$a \cdot \left<a\right>_0(b)$.

Prenant l'espace tangent \`a l'origine, on
d\'eduit de 5.11 un isomorphisme de pro-espaces
vectoriels de $\Lie\,V_\omega$ avec
$\scrL^{\uphat}$.
Transportons \`a $\scrL^{\uphat}$ l'action de
$\Lie\,V_\omega$ d\'eduite de l'action de
$V_\omega$ sur la structure 5.9.
D\'erivant (5.11.4), on trouve que 
$a\in \scrL^{\uphat}$ agit sur
$\scrL^{\uphat}=\Lie
\left(\prod\nolimits_{0,0}\right)$ par
$$
\partial_a\colon\, e_0\longmapsto 0,\quad
e_\zeta\longmapsto \left[-[\zeta](a),e_\zeta\right].
\tag5.12.2
$$
L'action sur $\prod\nolimits_{00}$ s'en
d\'eduit.
Il est clair sur (5.12.2) que
$\partial_{e_1}=0$.
Puisque le centralisateur de $e_1$ dans
$\scrL^{\uphat}$ est r\'eduit aux multiples de
$e_1$ (4.7), le noyau de $a\mapsto\partial_a$ est
r\'eduit \`a $\dbQ e_1$.

Pour \'ecrire l'action de
$\Lie\,V_\omega=\scrL^{\uphat}$ tant sur
$\prod\nolimits_{00}$ que sur le
$\prod\nolimits_{00}$-torseur
$\prod\nolimits_{1,0}$, il est commode de
plonger ces sch\'emas affines dans le dual de
leur alg\`ebre affine, un sch\'ema
provectoriel dont ils sont le sous-sch\'ema
des \'el\'ements groupaux.
A la structure de groupe de
$\prod\nolimits_{0,0}$ correspond une
structure d'alg\`ebre sur ce dual, \`a celle
de $\prod\nolimits_{0,0}$-torseur de
$\prod\nolimits_{0,1}$ une structure de module
\`a droite sous cette alg\`ebre:
$$
\spreadmatrixlines{2\jot}
\matrix
{\tsize\prod\nolimits_{0,1}}  &:\text{torseur
sous}\hfill &{\tsize\prod\nolimits_{0,0}} 
  &,\text{\'el\'ements groupaux de}\\
\null\kern-8pt\CAP &&\null\kern-8pt\CAP &\\ 
\dbQ\ll e_0,(e_\zeta)\gg 
   &:\text{module \`a droite sur}\hfill 
  &\dbQ\ll e_0,(e_\zeta)\gg &.
\endmatrix
$$

\medskip
Si $\dbQ[\eps]$ est l'alg\`ebre des
nombres duaux $(\eps^2=0)$, l'action de
$a\in\scrL^{\uphat}=\Lie(V_\omega)$ est le
coefficient de $\eps$ dans l'action du
$\dbQ[\eps]$-point $1+a\eps$ de $V_\omega$.
L'action sur l'alg\`ebre $\dbQ\ll
e_0,(e_\zeta)\gg$, duale de
$\scrO(\prod\nolimits_{0,0})$, est $e_0\mapsto
e_0$, $e_\zeta\mapsto
(1+[\zeta]a\eps)^{-1}e_\zeta(1+[\zeta]a\eps)$,
et correspond \`a l'action de a par la
d\'erivation $\partial_a\colon\, e_0\to 0$,
$e_\zeta\mapsto-[[\zeta]a,e_\zeta]$.
L'action sur le module \`a droite $\dbQ\ll
e_0,(e_\zeta)\gg$, dual de
$\scrO(\prod\nolimits_{0,1})$, est par
$x\mapsto (1+a\eps)(x+\partial_a x \cdot \eps)$ et
correspond \`a l'action de a par
$$
x\longmapsto ax+\partial_a x.\tag5.12.3
$$

Au crochet de Lie de $\Lie\,V_\omega$
correspond un nouveau crochet de Lie
$\{\,\,,\,\,\}$ sur $\scrL^{\uphat}$.

\proclaim{Proposition 5.13}
On a
$$
\{a,b\}=[a,b]+\partial_a(b)-\partial_b(a).
\tag5.13.1
$$
\endproclaim

\demo{Preuve}
V\'erifions d'abord que
$$
[\partial_a,\partial_b]=\partial_{[a,b]+\partial_a
(b)-\partial_b(a)}.\tag5.13.2
$$
Puisque, comme $\partial_a$ et $\partial_b$,
la d\'erivation $[\partial_a,\partial_b]$ de
$\scrL^{\uphat}$ fixe $e_0$ et est 
$\mu_N$-\'equivariante, il suffit de tester son action
sur $e_1$.
On a
$$
\align
\partial_a\partial_b &\colon\, e_1\longmapsto
  -[b,e_1]\longmapsto-[\partial_a(b),e_1]+
  [b,[a,e_1]]\\
\partial_b\partial_a &\colon\, e_1\longmapsto
  -[a,e_1]\longmapsto-[\partial_b(a),e_1]
  +[a,[b,e_1]]\\
[\partial_a,\partial_b] &\colon\,
e_1\longmapsto
-[\partial_a(b)-\partial_b(a)+[a,b],e_1].
\endalign
$$
\enddemo

Deux m\'ethodes pour d\'eduire (5.13.1) de
(5.13.2):

\medskip\noindent
{\sans a.}\enspace
Le crochet $\{\,\,,\,\,\}$ est compatible aux
graduations, toutes les constructions faites
\'etant compatibles \`a l'action de $\dbG_m$.
Le noyau de l'action de
$\scrL^{\uphat}=\Lie(V_\omega)$ sur $\scrL^{\uphat}$
\'etant r\'eduit aux multiples de $e_1$, si $a$
et  $b$ sont homog\`enes (de degr\'e $\Ge 1$),
il r\'esulte de (5.13.2) que les deux membres
de (5.13.1)
co{\umi}ncident \`a un multiple de $e_1$
pr\`es.
Etant homog\`ene de degr\'e $\Ge 2$, ils
co{\umi}ncident. 

\medskip\noindent
{\sans b.}\enspace
Testons sur (5.12.3).
Si $\mu_a$ est la multiplication \`a gauche
par $a$, on a
$$
\align
[\mu_a+\partial_a, &\mu_b+\partial_b] =
\mu_{[a,b]+\partial_ab-\partial_ba}+
  [\partial_a,\partial_b]\\
&=\mu_{[a,b]+\partial_ab-\partial_b a}+
\partial_{[a,b]+\partial_ab-\partial_ba}.
\endalign
$$

\remark{\bf 5.14 Remarque}\enspace
(i) On a
$$
\{a,e_1\}=[a,e_1]+\partial_a(e_1)-\partial_{e_1}(a)
=[a,e_1]-[a,e_1]+0=0:
$$
$e_1$ est dans le centre de
$(\scrL^{\uphat},\{\,\,,\,\,\})$.

\medskip\noindent
(ii)
On prendra garde que l'application
exponentielle de Lie $V_\omega$ dans
$V_\omega$, identifi\'ee par 5.11 et 5.12 \`a
une application de $\scrL^{\uphat}$ dans $\prod$,
n'est pas l'application exponentielle de
$\scrL^{\uphat}$ dans $\prod$.
Il resulte de 5.12.4 que, si $\mu_a$ est la
multiplication \`a gauche par $a$, c'est
$$
\align
\exp &\colon\, \Lie\,V_\omega\to V_\omega
  \colon\,\scrL^{\uphat}\to\prod\Subset\dbQ
\ll e_0,(e_\zeta)\gg\colon\\
&a\longmapsto\sum (\mu_a+\partial_a)^n/n!
\quad\text{appliqu\'e \`a $1$}.
\endalign
$$
Premiers termes:
$$
a+(a^2+\partial_a(a))/2+(a^3+2a \cdot \partial_a(a)+
\partial_a(a) \cdot a+\partial_a^2(a))/
6+\cdots
$$
\endremark

\remark{\bf Remarque 5.15}
Ainsi que la notation le sugg\`ere, le groupe
$H_\omega$ (resp. $V_\omega$) est la
r\'ealisation $\omega$ d'un
$\MT(\scrO[1/N])$-sch\'ema en groupe $H$ (resp. $V$).
On peut d\'efinir $H$ par la propri\'et\'e que
pour tout foncteur fibre $F$, $F(H)$ est,
fonctoriellement en $F$, le sch\'ema en
groupe des automorphismes de la structure 5.7
correspondante.
Il agit sur $\dbQ(1)$, $P_{0,0}$, et $P_{0,1}$
et $V$ est le noyau de l'homomorphisme
$H\to\dbQ(1)$ donnant l'action de $H$ sur
$\dbQ(1)$.
Il reste vrai dans $\MT(\scrO[1/N])$ que $P_{1,0}$
est un espace principal homog\`ene sous $V$.
\endremark

Comme en 5.10, le groupe
$G:=\pi(\MT(\scrO[1/N]))$ s'envoie dans $H$, et
ce morphisme induit un morphisme
$$
U\to V\tag5.15.1
$$
de $\MT(\scrO[1/N])$-sch\'emas en groupes.
Sa r\'ealisation $\omega$ est induite par
(5.10.3).

\subhead
5.16
\endsubhead
Fixons un plongement $\sigma$ de $k$ dans
$\dbC$, i.e. un isomorphisme de $\mu_N(k)$
avec $\mu_N(\dbC)$.
Le {\it droit chemin} du point base $1_0$ de
$\dbC^*-\mu_N$ vers le point base $(-1)_1$ va
par segments de droite successifs \quad
(a) dans l'espace tangent en $0$, de $1$ \`a
$\eps>0$; (b) de $\eps$ \`a $1-\eps$ dans
$\dbC^*-\mu_N$; (c) de $-\eps$ \`a $-1$ dans
l'espace tangent \`a un.
Ce droit chemin a une image $\dch$ dans
$(P_{(-1)_1,1_0})_\sigma$.
Via l'isomorphisme de comparaison
$\comp_{\sigma,\omega}$, il correspond \`a un
\'el\'ement de 
$$
\dch(\sigma)\in (P_{(-1)_1,1_0})_\omega\otimes\dbC
=\prod(\dbC) \Subset \dbC\ll
e_0,(e_\zeta)_{\zeta\in\mu_N(k)}\gg.
$$

Les coefficients de $\dch(\sigma)$ sont
donn\'es par 
$$
\spreadlines{1\jot}
\leqalignno{
&\text{coefficient de }e_{x_1}\ldots 
     e_{x_n}= &(5.16.1)\cr
&=\text{ int\'egrale it\'er\'ee de $0$ \`a $1$ de }
\frac{dz}{z-\sigma(x_1)}\cdots
\frac{dz}{z-\sigma(x_n)}\cr
&=\int_0^1\left(\frac{dz_1}{z_1-\sigma(x_1)}
\int_0^{z_1}\left(\frac{dz_2}{z_2-\sigma(x_2)}
\cdots \int_0^{z_{n-1}}
\frac{dz_n}{z_n-\sigma(x_n)}\right)\ldots
\right).\cr}
$$
Si $x_n=0$ ou si $x_1=1$, l'int\'egrale diverge
et doit \^{e}tre r\'egularis\'ee:
remplacer les limites d'int\'egration $0$ et
$1$  par $\eps$ et $1-\eta$; l'int\'egrale
it\'er\'ee convergente obtenue a pour comportement
asymptotique pour $\eps,\eta\to 0$
$$
\align
\text{Int\'egrale it\'er\'ee}
&=\text{polyn\^{o}me
$(\log\,\eps,\log\,\eta)$}\\
&\qquad +O(\sup(\eps\vert\log\,\eps\vert^A,
\eta\vert\log\,\eta\vert^B))
\endalign
$$
pour $A$ and $B$ convenables, et on prend le
terme constant du polyn\^{o}me.

Les valeurs multizeta et leur analogue
cyclotomique apparaissent quand dans (5.16.1)
on d\'eveloppe $dz/z-\zeta$ en s\'erie
g\'eom\'etrique:
$$
\frac{dz}{z-\zeta}=\frac{-\zeta^{-1}z}{1-\zeta^{-1}z}
\,\,\frac{dz}{z}=-\sum\limits_{n\Ge 1}(\zeta^{-1}z)^n.
\frac{dz}{z}.
$$

\proclaim{Proposition 5.17}
{\rm (i)}\enspace Pour $\zeta_1,\dotsc,\zeta_m$ dans
$\mu_N(k)$ et des entiers $s_1,\dotsc,s_m\Ge
1$, si $s_1\not=1$, le coefficient de
$e_0^{s_1-1}e_{\zeta_1}\ldots
e_0^{s_m-1}e_{\zeta_m}$ dans $\dch(\sigma)$ est
$$
(-1)^m\sum\limits_{n_1>\ldots>n_m>0}
\frac{\sigma\left(\zeta_1^{n_2-n_1}\zeta_2^{n_3-n_2}
\ldots \zeta_m^{-n_m}\right)}
{n_1^{s_1}\ldots n_m^{s_m}}
$$

\noindent
{\rm (ii)}\enspace
$\dch(\sigma)$ est caract\'eris\'e par les
propri\'et\'es d'\^{e}tre un \'el\'ement
groupal (terme constant 1, et 
$\Delta\dch(\sigma)=\dch(\sigma)\otimes
\dch(\sigma)$) ayant les coefficients (i) et
dont les coefficients de $e_0$ et $e_1$ sont
nuls.
\endproclaim

Notons $\tau$ l'application de $\dbG_m$ dans
$H_\omega$ d\'efinie par les graduations
(5.10.2) et $\left<g\right>$ l'\'el\'ement de
$V_\omega(\dbC)$ correspondant par 5.11 \`a
$g\in\prod(\dbC)$.

\proclaim{Proposition 5.18}
L'\'el\'ement
$\left<\dch(\sigma)\right>\tau(2\pi i)$
de $H_\omega(\dbC)$ transforme les 
$\dbQ$-formes $(P_{0,0})_\omega$ et
$(P_{1,0})_\omega$ de
$(P_{0,0})_\omega\otimes\dbC$ et
$(P_{1,0})_\omega\otimes\dbC$ en les
$\dbQ$-formes $(P_{0,0})_\sigma$ et
$(P_{1,0})_\sigma$.
\endproclaim

\demo{Preuve}
Les \'el\'ements de $\pi_1(X(\dbC),1_0)$ ont
dans $\prod(\dbC)$ une image rationnelle pour
la $\dbQ$-forme $(P_{0,0})_\sigma$.
Ce $\pi_1$ contient la monodromie locale en
$0$ et le chemin allant droit pr\`es de $1$,
tournant autour de $1$, et revenant.
Images: $\exp(2\pi i\,e_0)$ et
$\dch(\sigma)^{-1}\exp(2\pi
i\,e_1)\dch(\sigma)$.
Ce sont les images de $\exp(e_0)$ et
$\exp(e_1)$ par
$\left<\dch(\sigma)\right)\tau(2\pi i)$.
La $\dbQ$-forme $(P_{0,0})_\sigma$ est
$\mu_N$-\'equivariante.
Les images des $\exp(e_\zeta)$ par
$\left<\dch(\sigma)\right>\tau(2\pi i)$ sont
donc elles aussi rationnelles.
La pro-alg\`ebre de Lie $\scrL^{\uphat}$ est
l'unique $\dbQ$-forme de
$\scrL_{\dbC}^{\uphat}$ pour laquelle les
$e_0$ et $e_\zeta$ sont rationnels, et $\prod$
est donc l'unique $\dbQ$-forme du groupe
$\prod_{\dbC}$ pour laquelle $\exp(e_0)$ et
les $\exp(e_\zeta)$ sont rationnels.
Transportant par
$\left<\dch(\sigma)\right>\tau(2\pi i)$, on
obtient l'assertion pour $(P_{0,0})_\omega$.

Montrons que la $\dbQ$-forme
$(P_{1,0})_\sigma$ de
$(P_{1,0})_\omega\otimes\dbC$ est l'image de
la $\dbQ$-forme $(P_{1,0})_\omega$ par
$\left<\dch(\sigma)\right>\tau(2\pi i)$.
Ces deux $\dbQ$-formes sont des
$(P_{0,0})_\sigma$-torseurs.
Il suffit donc de montrer qu'elles admettent
un point rationnel commun.
L'image $\dch(\sigma)$ de $1$ par
$\left<\dch(\sigma)\right>\tau(2\pi i)$ est en
effet dans $(P_{1,0})_\omega(\dbQ)$ car image
du droit chemin de $1_0$ \`a $(-1)_1$.
\enddemo

\subhead
5.19
\endsubhead
Notons encore $a_\sigma$ un \'el\'ement de
$G_\omega\left<\MT(\scrO[1/N])\right>$
image d'un \'el\'ement\break
 $a_\sigma\in G_\omega\left<\MT(k)\right>$ de 2.12, et
$\iota(a_\sigma)$ son image par 5.10.3 dans
$H_\omega$.
De m\^{e}me pour $a_\sigma^0$ (2.12.2).
L'\'el\'ement $\iota(a_\sigma)$ lui-aussi
transforme les $\dbQ$-formes $\omega$ de
$(P_{00})_\omega\otimes\dbC$ et
$(P_{1,0})_\omega\otimes\dbC$ en les
$\dbQ$-formes $\sigma$.
Il existe donc $v_\sigma\in H_\omega(\dbQ)$
tel que
$$
\left<\dch(\sigma)\right>\tau(2\pi
i)=\iota(a_\sigma)v_\sigma.\tag5.19.1
$$

Si comme en (2.12.2) on a choisit $a_\sigma$
de la forme $a_\sigma^0\tau(2\pi i)$ avec
$a_\sigma^0\in U_\omega(\dbC)$, tant
$\left<\dch(\sigma)\right>\tau(2\pi i)$ que
$\iota(a_\sigma)$ ont la m\^{e}me projection
$2\pi i$ dans le quotient $\dbG_m$ de
$H_\omega$, $v_\sigma$ est dans
$V_\omega(\dbQ)$ et on peut r\'ecrire (5.19.1)
sous les formes
$$
\gather
\left<\dch(\sigma)\right>=\iota(a_\sigma^0)\cdot
  (\tau(2\pi i)v_\sigma\tau(2\pi i)^{-1})\tag5.19.2\\
\tau(2\pi i)^{-1}\left<\dch(\sigma)\right>
   \tau(2\pi i)=
  (\tau(2\pi i)^{-1}\iota(a_\sigma^0)
  \tau(2\pi i))v_\sigma.\tag5.19.3
\endgather
$$

A ces formules correspondent deux fa\c{c}ons
de voir les relations entre les coefficients
(5.17.1) de $\dch(\sigma)$ impliqu\'ees par la
th\'eorie motivique.

La formule (5.19.2) implique que
$\dch(\sigma)\in\prod(\dbC)$ est contenu dans
la sous-vari\'et\'e alg\'ebrique d\'efinie sur
$\dbQ$ suivante.
Soit $\iota(U_\omega)\Subset\left(\prod,\circ\right)$
le sous-groupe unipotent image de $U_\omega$
par (5.10.3).
Soit $D_\sigma$ l'image de la droite de
coordonn\'ee $t$ par $t\mapsto
\tau(t)v_\sigma\tau(t)^{-1}$.
Cette application est d\'efinie en $t=0$ car
$\scrL^{\uphat}$ est \`a degr\'es $>0$.
On a $t=0\mapsto$ \'el\'ement neutre.
Par (5.19.2), $\dch(\sigma)$ est contenu dans
l'image par le produit $\circ$ de
$\iota(U_\omega)$ et $D_\sigma$.
Noter que $D_\sigma$ peut d\'ependre de
$\sigma$, mais $\iota(U_\omega)$ non.

La formule (5.19.3) sugg\`ere de consid\'erer
plut\^{o}t $\tau(2\pi i)^{-1}(\dch(\sigma))$,
de coefficients convergents donn\'es par la
formule 5.17 (i), divis\'ee par $(2\pi
i)^{\sum s_i}$.
Les autres coefficients sont d\'eduits de
ceux-l\`a par des formules rationnelles.
L'\'el\'ement $\tau(2\pi
i)^{-1}(\dch(\sigma))\in\prod(\dbC)$ est
contenue dans la sous-vari\'et\'e d\'efinie
sur $\dbQ$ $\iota(U_\omega)v_\sigma$, une
classe lat\'erale rationnelle pouvant
d\'ependre de $\sigma$ du sous-groupe
$\iota(U_\omega)$ de $\left(\prod, \circ \right)$

\subhead
5.20
\endsubhead
Supposons que $N=1$.
On a $k=\dbQ$; $\sigma$ est unique.
Choisissons $a_\sigma^0$ r\'eel (2.12).
Puisque $\dch(\sigma)$ est r\'eel, (5.12.2)
montre que $\tau(2\pi i)v_\sigma\tau(2\pi
i)^{-1}$ est r\'eel.
Le logarithme de $v_\sigma$ est donc purement
en degr\'e pair, $\tau(t)v_\sigma\tau(t)^{-1}$
ne d\'epend que de $t^2$ et ceci permet de
poser $d(t^2):=\tau(t)v_\sigma\tau(t)^{-1}$.

Regardons $\prod$ comme l'espace des
\'el\'ements groupaux de $\dbC\ll e_0,e_1\gg$.
L'alg\`ebre de Lie de $U_\omega$ \'etant en
degr\'e $\Ge 3$, les coefficients de
$\dch(\sigma)$ et de $v_\sigma$
co{\umi}ncident en degr\'ee $\Le 2$.
On a 
$$
\dch(\sigma)=1-\zeta(2)[e_0,e_1]+\text{degr\'e}\Ge 3
$$
et donc, puisque $\zeta(2)=-(2\pi i)^2/24$,
$$
v_\sigma=1+\tfrac{1}{24}[e_0,e_1]+\text{degr\'e}\Ge 3.
$$

Pour $g\in\iota(U_\omega(\dbC))$ et $d\in
D_\sigma$, $d=\tau(t)v_\sigma\tau(t)^{-1}$,
$t^2$ est donc $24$ fois le coefficient de
$e_0e_1$ dans $d$, \'egal \`a celui de $e_0e_1$
dans $g\circ d$.
Ceci d\'etermine $d$ \`a partir de $g\circ d$
et
$$
\iota(U_\omega)\times D_\sigma\to
\left(\prod, \circ \right)
$$
est un plongement ferm\'e: l'image est
l'ensemble des $x\in\prod\Subset\dbC\ll
e_0,e_1\gg$ tels que
$$
x\circ d\text{ (24.coefficient de $e_0e_1$
dans $x$)}^{-1}\in\iota(U_\omega).
\tag5.20.1
$$

Une variante d'une conjecture de Grothendieck
affirme que $a_\sigma$ est Zariski-dense (sur
$\dbQ$) dans $G_\omega$.
Si tel est le cas, toutes les relations
alg\'ebriques entre les coefficients de
$\dch_\sigma$ sont donn\'ees par 5.20.1.

On observera que les \'equations polynomiales
entre les coefficients de $\dch_\sigma$
d\'eduites de 5.20.1 ont les propri\'et\'es
suivantes.

\medskip\noindent
{\sans (a)}\enspace
Homog\'en\'eit\'e: toute \'equation est somme
d'\'equations isobares, pour le poids.
Ceci exprime que $\iota(U_\omega)\circ
D_\sigma\Subset V_\omega$ est stable sous
l'action int\'erieure de $\dbG_m\Subset
H_\omega$.
En effet, tant $\iota(U_\omega)$ que
$D_\sigma$ le sont.

\smallskip\noindent
{\sans (b)}\enspace
Si on ajoute l'\'equation ``$\pi^2=0$'',
plus pr\'ecis\'ement: ``nullit\'e du
coefficient de $e_0e_1$'', on obtient l'id\'eal 
des \'equations d\'efinissant le
sous-groupe $\iota(U_\omega)$ de
$\left(\prod,0\right)$.

\smallskip\noindent
{\sans (c)}\enspace
Si on ajoute plut\^{o}t l'\'equation
``$\pi^2=\alpha$'', on obtient une classe
lat\'erale de ce groupe.

\medskip
G. Racinet donne dans sa th\`ese un syst\`eme
d'\'equations polynomiales v\'erifi\'ees par
les coefficients de $\dch(\sigma)$ et prouve
que ce syst\`eme v\'erifie (a) et la forme
affaiblie suivante de (b) (c): si on ajoute
l'\'equation ``$\pi^2=0$'', on obtient un
sous-groupe de $\left(\prod,0\right)$, et si
on ajoute ``$\pi^2=\alpha$'', on obtient une
classe lat\'erale de ce sous-groupe.

\subhead
5.21
\endsubhead
La situation est similaire pour $N=2$,
$\zeta(2)$ continuant \`a jouer le m\^{e}me
r\^{o}le que pour $N=1$.
Il n'est plus vrai que $\Lie\,U_\omega$ est en
degr\'e $\Ge 3$, mais on est sauv\'e par la
compatibilit\'e entre les cas $N=1$ et $N=2$:
$$
\spreadlines{2\jot}
\CD
U_\omega\left<\MT(\dbZ[1/2])\right> @>>>
\left(\prod\text{ pour }N=2,o\right)\\
@VVV @VVV\\
U_\omega\left<\MT(\dbZ)\right> @>>>
  \left(\prod\text{ pour }N=1,o\right).
\endCD
$$

\subhead
5.22
\endsubhead
Pour $N\Ge 3$, $k$ est totalement complexe.
Le coefficient de $e_\zeta$ dans
$\dch(\sigma)$ est $-\log(1-\sigma(\zeta))$
(resp. $0$ pour $\zeta=1$), et le r\^{o}le
pr\'ec\'edemment jou\'e par le coefficient
$-\zeta(2)$ de $e_0e_1$ est jou\'e par la
diff\'erence des coefficients de $e_\zeta$
et de $e_{\zeta^{-1}}$, pour $\zeta\in\mu_N$
outre que $\pm1$.
Cette diff\'erence est un multiple rationnel
de $2\pi i$, d\'ependant de $\sigma$.
La courbe $D_\sigma$ est encore une droite,
param\'etr\'ee par la difference
$\delta_{\coeff}$ des coefficients de
$e_\zeta$ et $e_{\zeta^{-1}}$, 
$$
\iota(U_\omega)\times D_\sigma\to
\left(\prod, \circ\right)
$$
est un plongement ferm\'e, et l'image est
l'ensemble des $x\in\prod\Subset\dbC\ll
e_0,(e_\alpha)_{\alpha\in\mu_N}\gg$ tels que
$$
x\circ\text{(point de $D_\sigma$ donn\'e par
$\delta_{\coeff}(x))^{-1}\in\iota(U_\omega)$}.
\tag5.22.1
$$

La conjecture de Grothendieck affirme \`a
nouveau que toutes les relations polynomiales
\`a coefficients rationnels entre les
coefficients de $\dch(\sigma)$ proviennent de
(5.22.1).

Le syst\`eme d'\'equations polynomiales
d\'eduit de (5.22.1) a encore les
propri\'et\'es (a), (b), (c) de 5.20, avec
``$\pi^2$'' remplac\'e par ``$2\pi i\,$'' (plus
pr\'ecis\'ement, par la diff\'erence des
coefficients de $e_\zeta$ et
$e_{\zeta^{-1}}$), et Racinet a obtenu 
pour tout $N$ des r\'esultats semblables \`a ceux
mentionn\'es en 5.20.

\subhead
5.23
\endsubhead
Permettons \`a nouveau \`a $N$ d'\^{e}tre un
quelconque entier $\Ge 1$.
Puisque $\dch(\sigma)$ est un \'el\'ement
groupal de
$\dbC\ll e_0,(e_\alpha)_{\alpha\in\mu_N}\gg$, le
produit de deux coefficients de $\dch(\sigma)$
est combinaison lin\'eaire \`a coefficients
rationnels de coefficients, de sorte que
conna{\hati}tre les relations lin\'eaires entre les
coefficients suffit pour conna{\hati}tre  les
relations polynomiales entre eux.

\proclaim{Th\'eor\`eme 5.24}
Soit $d_n$ la dimension en degr\'e $n$ de
l'image de $\iota\colon\,\Lie\,U_\omega\to
\Lie\,V_\omega$ (5.10.3).
Soit $A$ l'alg\`ebre de polyn\^{o}mes
gradu\'ee engendr\'ee sur $\dbQ$ par $d_n$
g\'en\'erateurs degr\'e $n$, et par un
g\'en\'erateur additionnel $t_0$ en degr\'e
$1$ (resp. $2$ si $N=1$ ou $2$).

Il existe un homomorphisme $\varphi_\sigma$ de
$A$ dans $\dbC$, envoyant $t_0$ sur $2\pi i$
(resp. $\pi^2$ si $N=1$ ou $2$), tel que les
coefficients de mon\^{o}mes de degr\'e $d$
dans $\dch_\sigma$ soient contenus
dans\break
 $\varphi_\sigma$ (partie de degr\'e $d$ de $A$).
\endproclaim

\demo{Preuve}
On prend simplement pour $A$ l'alg\`ebre
affine de $\iota(U_\omega)\times D_\sigma$.
Elle a la structure dite d'alg\`ebre
gradu\'ee, la fonction ``coefficient d'un
mon\^{o}me de degr\'e $d$'' est de degr\'e
$d$, et $\dch(\sigma)$ est un point de
$\Spec(A)$ \`a valeurs dans $\dbC$, i.e. un
homomorphisme $\varphi_\sigma$ de $A$ dans
$\dbC$.
Le g\'en\'erateur $t_0$ est la coordonn\'ee de
$D_\sigma$, convenablement normalis\'ee.
\enddemo

\proclaim{Corollaire 5.25}
Soit $D_n$ la dimension du $\dbQ$-espace
vectoriel engendr\'e par les coefficients des
monomes de degr\'e $n$ dans $\dch(\sigma)$.
Le s\'erie g\'en\'eratrice $\sum D_nt^n$ est
terme \`a terme major\'ee par la suivante:

\roster
\item"{\sans (i)}"
Pour $N=1$: $1/(1-t^2-t^3)$.

\item"{\sans (ii)}"
$N=2$: $1/(1-t-t^2)$.

\item"{\sans (iii)}"
Pour $N\Ge 3$ ayant $\nu$ facteurs premiers:
$1/\left(1-\left(\frac{\varphi(N)}{2}+(\nu-1)\right)
t+(\nu-1)t^2\right)$, o\`u $\varphi$
est l'indicateur d'Euler
$\vert(\dbZ/N\dbZ)^*\vert$.
\endroster
\endproclaim

\demo{Preuve}
La s\'erie majorante est la s\'erie
g\'en\'eratrice des dimensions des composantes
homog\`enes de l'alg\`ebre gradu\'ee de
polyn\^{o}me construite comme l'alg\`ebre $A$
de 5.25, mais en partant de $\Lie(U_\omega)$
plut\^{o}t que d'un quotient.
En effet, cette s\'erie est le produit de
$\frac{1}{1-t}$ (resp $\frac{1}{1-t^2}$ si
$N=1$ ou $2$) par la s\'erie g\'en\'eratrice
pour l'alg\`ebre enveloppante de
$\Lie(U)_\omega^{\GR}$.
Cette derni\`ere est une alg\`ebre associative
libre.
S\'erie g\'en\'eratrice: $\sum
f(t)^n=1/(1-f(t))$, avec
$$
f(t)=\frac{\varphi(N)}{2}\,\,\frac{t}{1-t}+(\nu-1)t
$$
(resp. $f(t)=t^3/1-t^2$ si $N=1$, $t/1-t^2$ si
$N=2$).
Le calcul final est laiss\'e au lecteur.
\enddemo

\remark\nofrills{\bf Remarque 5.26.}\enspace
Pour $N=1$, nous retrouvons ainsi le
r\'esultat de Terasoma (2002).
Terasoma consid\`ere un groupe de cohomologie
relative, correspondant \`a un motif de Tate
mixte sur $\dbZ$ qui est sans doute le
suivant: dans l'alg\`ebre affine de ${}_1P_0$,
prendre la partie de poids $\Le 2n$, puis les
coinvariants de la monodromie locale en $0$
et $1$, agissant \`a droite (resp. gauche) sur
$P_{1,0}$.
Comme la n\^{o}tre, sa preuve repose sur 1.6.
\endremark

\subhead
5.27
\endsubhead
La conjonction de (a) toute relation
$\dbQ$-lin\'eaire entre les coefficients de
$\dch(\sigma)$ est somme de relations isobares
(b) la borne 5.25 est atteinte \'equivaut \`a
la conjonction de (c) la conjecture de
Grothendieck que $a_\sigma$ est $\dbQ$-Zariski
dense dans $U_\omega$, et (d) l'injectivit\'e
de $\iota\colon\, U_\omega\to V_\omega$.
La condition (d) est en g\'en\'eral fausse.
Elle est vraie pour $N=2$, $3$ ou $4$ et on
ignore si elle est vraie pour $N=1$.

Pour obtenir une estimation raisonable, il
faudrait conna{\hati}tre la dimension des
composantes gradu\'ees l'image de $\Lie\,U_\omega$
dans $\Lie\,V_\omega$.
On ne conna{\hati}t ces dimensions que pour
$N=2,3,4$.
La question plus int\'eressante serait de
d\'eterminer non seulement ces dimensions, mais
encore l'image elle-m\^{e}me de $\Lie\,U_\omega$.
On ne sait le faire dans aucun cas.

\subhead
5.28
\endsubhead
Il serait tr\`es int\'eressant de disposer
pour les motifs de Tate mixtes de
r\'ealisations cristallines qui, appliqu\'ees
aux espaces de chemins motiviques, fournissent
les r\'ealisations consid\'er\'ees en Deligne
(1989).

Dans le cas de motifs de Tate mixtes $M$ sur
$\dbZ[1/D]$, on veut, pour chaque nombre
premier $p$ premier \`a $D$, un automorphisme
dit de Frobenius
$$
F_p\colon\, M_\sigma\otimes\dbQ_p\to
M_\sigma\otimes\dbQ_p.
$$
Cet automorphisme est l'action d'un
\'el\'ement canonique de $G_\omega(\dbQ_p)$,
pour
$$
G_\omega=G_\omega\left<\MT(\dbZ[1/D])\right>.
$$
Sur $\dbQ(1)_\sigma\otimes\dbQ_p$, $F_p$ est
la multiplication par $1/p$.

Pour $N=1$ ou $2$, on a $k=\dbQ$ et, pour $p$
premier \`a $N$, on disposerait donc de
$F_p^{-1}$ de la forme
$$
F_p^{-1}=\varphi_p\tau(p)\in G_\omega(\dbQ_p)
$$
et de l'image $\iota(\varphi_p)$ de
$\varphi_p$ dans $V_\omega(\dbQ_p)$.
Plongeons comme d'habitude $V_\omega(\dbQ_p)$
dans\break
$\dbQ_p\ll e_0,e_1\gg$, pour $N=1$, ou
$\dbQ_p\ll e_0,e_1,e_{-1}\gg$, pour $N=2$.
Les coefficients de $\iota(\varphi_p)$
semblent \^{e}tre des analogues $p$-adiques des
valeurs multiz\^{e}ta (cf Deligne (1989) 3.4).
D'apr\`es 5.20, on s'attend \`a ce qu'ils
v\'erifient toutes les identit\'es satisfaites
par les coefficients de $\dch(\sigma)$, ainsi
que ``$\pi^2=0$'' (nullit\'e du coefficient de
$e_0e_1$).
Le formalisme cristallin esp\'er\'e prouverait
en tout cas les identit\'es d\'eduites de
(5.20.1), et de ``$\pi^2=0$''.

Il serait int\'eressant aussi de disposer pour
ces coefficients d'expressions $p$-adiques qui
rendent claires qu'ils v\'erifient des
identit\'es du type
$$
\align
&\coeff(e_0^{n-1}e_1)\coeff(e_0^{m-1}e_1)=\\
&\coeff(e_0^{n-1}e_1e_0^{m-1}e_1)+\coeff
  (e_0^{m-1}e_1e_0^{n-1})\\
&\qquad\qquad\qquad +\coeff(e_0^{n+m-1})e_1),
\endalign
$$
qui pour $\dch(\sigma)$ expriment que
$$
\sum \frac{1}{k^n}.\sum \frac{1}{\ell^m}=
\sum\limits_{k>\ell}\frac{1}{k^n}
\frac{1}{\ell^m}+\sum\limits_{\ell>k}
\frac{1}{k^n}\frac{1}{\ell^m}+\sum
\frac{1}{k^{n+m}}.
$$

Pour $N\Ge 3$, le plus commode est de se
rappeller que, m\^{e}me si les points de
$\mu_N$ ne sont d\'efinis que sur $k$,
$X=A^*-\mu_N$ est d\'efini sur $\dbQ$, de
sorte que $\pi_1(X,0)$ et $P_{1,0}(X)$ sont de
fa\c{c}on naturelle des
$\MAT(k/\dbQ)$-sch\'emas.
La r\'ealisation de de Rham de
$\Lie\,\pi_1(X,0)$ est une $\dbQ$-forme
$\scrL_{\DR}^{\uphat}$ de
$\scrL^{\uphat}\widehat{\otimes} k$: les
invariants de $\Gal(k/\dbQ)$ agissant sur $k$
et permutant les $e_\zeta$ ($\zeta\in\mu_N$).
Pour $p\nmid N$, la th\'eorie cristalline
fournit un automorphisme de
$\scrL_{\DR}^{\uphat}$ et du
$\exp(\scrL_{\DR}^{\uphat})$-torseur
$P_{1,0}(X)_{\DR}$.




